% Current evidence-safe SCD-ICML results fragment.
% Generated on 2026-06-10 from verified MistralLite -> Mistral-7B runs.
% This fragment is intended to replace the historical broad result tables in
% the archived ICML camera-ready source if a current reproducibility-safe paper
% revision is needed. Qwen-32B is intentionally deferred.

\subsection{Current Reproducible Result Boundary}
\label{subsec:current_repro_boundary}

We report the current machine-verified SCD result on the MistralLite$\to$Mistral-7B
pair.  The release path uses a fixed no-fallback, no-recompute route: 19 layers
are served by \textsc{Reuse}, and 13 transition layers are served by
\textsc{Patch}.  The route is selected by F1/EM retention on CMRC2018 and is
stored in the release package as
\texttt{configs/routes/scd\_icml\_mistrallite\_to\_mistral7\_reuse19\_patch13\_cmrc200.json}.
Qwen-32B is not included in this release-safety audit.

\begin{table}[t]
  \centering
  \small
  \caption{\textbf{Verified CMRC2018 selected-route result.}
  MistralLite$\to$Mistral-7B, CMRC2018 dev, $n=200$, greedy decoding, no
  fallback, no recomputation.  SCD preserves F1/EM relative to Full Consumer
  while reducing TTFT.  Avg NLL/PPL are reported as auxiliary calibration-risk
  signals, not as the route-selection objective.}
  \label{tab:current_cmrc_selected_route}
  \setlength{\tabcolsep}{4pt}
  \begin{tabular}{lrrrrr}
    \toprule
    \textbf{Method} & \textbf{F1} & \textbf{EM} & \textbf{Avg NLL} & \textbf{PPL} & \textbf{TTFT (ms)} \\
    \midrule
    Full Consumer & 0.657307 & 0.365 & 0.569284 & 1.767002 & 100.604 \\
    SCD Reuse19/Patch13 & 0.642751 & 0.350 & 0.934572 & 2.546122 & 41.598 \\
    \bottomrule
  \end{tabular}
\end{table}

The verified F1 retention is $0.977855$ and EM retention is $0.958904$.
Thus, under the F1/EM-primary objective, the current release result does not
show evidence of fabrication or an unrunnable main path: the open-source wrapper
reaches the selected route, and the completed remote run reproduces the same
Full Consumer/SCD comparison.

\subsection{Maximum-Reuse Boundary}
\label{subsec:current_reuse_boundary}

The route-selection objective is to maximize the number of \textsc{Reuse} layers
while preserving answer quality.  Starting from the Reuse19/Patch13 route, each
single-layer extension to Reuse20 was tested.  The best Reuse20 candidate
already falls below the working F1-retention target, and higher-Reuse routes
degrade rapidly.

\begin{table}[t]
  \centering
  \small
  \caption{\textbf{Quality boundary as Reuse count increases.}
  CMRC2018 dev, $n=200$.  Reuse19/Patch13 is the current maximum verified
  quality-preserving route under the F1/EM-primary criterion.}
  \label{tab:current_reuse_boundary}
  \setlength{\tabcolsep}{5pt}
  \begin{tabular}{lrrr}
    \toprule
    \textbf{Route} & \textbf{CMRC F1} & \textbf{CMRC EM} & \textbf{Interpretation} \\
    \midrule
    Reuse19/Patch13 & 0.642751 & 0.350 & current boundary \\
    Reuse20/Patch12 & 0.611265 & 0.360 & below F1-retention target \\
    Reuse21/Patch11 & 0.554634 & 0.310 & too much quality loss \\
    Reuse22/Patch10 & 0.496437 & 0.255 & too much quality loss \\
    Reuse29/Patch03 & 0.077045 & 0.010 & collapsed \\
    Reuse30/Patch02 & 0.009595 & 0.000 & collapsed \\
    Reuse31/Patch01 & 0.007903 & 0.000 & collapsed \\
    \bottomrule
  \end{tabular}
\end{table}

\subsection{Limitations on Other Benchmarks}
\label{subsec:current_limitations}

The current selected route should not be presented as a universal
three-dataset quality-preserving result.  On full-context HotpotQA, the
CMRC-selected route loses substantial F1/EM even after retraining HotpotQA banks
on a disjoint split.  On WikiText-2, the selected route improves TTFT but worsens
language-modeling quality as measured by NLL/PPL.  These rows are therefore
limitations and diagnostics, not positive main claims.

\begin{table}[t]
  \centering
  \small
  \caption{\textbf{Current non-CMRC diagnostics.}
  HotpotQA uses complete distractor context with no truncation
  (\texttt{MAX\_LENGTH=0}).  WikiText-2 reports teacher-forced
  language-modeling quality.}
  \label{tab:current_non_cmrc_diagnostics}
  \setlength{\tabcolsep}{4pt}
  \begin{tabular}{llrrr}
    \toprule
    \textbf{Dataset} & \textbf{Method} & \textbf{F1/EM or NLL/PPL} & \textbf{TTFT (ms)} & \textbf{Conclusion} \\
    \midrule
    HotpotQA & Full Consumer & F1 0.400453 / EM 0.265 & 216.153 & baseline \\
    HotpotQA & SCD Reuse19/Patch13 & F1 0.138272 / EM 0.105 & 67.060 & quality not preserved \\
    WikiText-2 & Full Consumer & NLL 2.026045 / PPL 7.584029 & 62.368 & baseline \\
    WikiText-2 & SCD Reuse19/Patch13 & NLL 3.274239 / PPL 26.423097 & 31.004 & auxiliary risk \\
    \bottomrule
  \end{tabular}
\end{table}

An older GitHub snapshot reports a stronger HotpotQA Patch-only result
(F1 $0.351659$, EM $0.252127$ on 7405 examples), but that evaluation uses a
question-aware compressed-evidence prompt and \texttt{max\_length=512}.  It is a
separate compressed-prompt benchmark and should not be mixed with the current
full-context HotpotQA setting; the release code exposes it as
\texttt{hotpotqa\_compressed} for isolated reproduction.  A question-text overlap audit found that this
old 7405-example evaluation includes the 100 questions used to train the old
Patch checkpoint; after removing those questions, Patch-only remains at F1
$0.351443$ and EM $0.251608$.  Thus the result is real for the compressed
setting, but it is not evidence for full-context HotpotQA.

\paragraph{Evidence roots.}
The CMRC selected-route aggregate is
\texttt{/disk\_n/c101-3-mql/projects/SCD-ICML/runs/scd\_icml\_selected\_reuse19\_5method\_aggregate\_20260610\_034426}.
The valid full-context HotpotQA run is
\texttt{/disk\_n/c101-3-mql/projects/SCD-ICML/runs/hotpotqa\_fullcontext\_full\_scd\_20260610\_115952}.
The HotpotQA full-bank heldout run is
\texttt{/disk\_n/c101-3-mql/projects/SCD-ICML/runs/hotpotqa\_full\_banks\_20260610\_192156}.
